\chapter{Preparing a Station: \texttt{ate setup}}\label{ch:setup}

One command prepares a Station however ATE was installed. It is idempotent: run it again and it
creates nothing that already exists and never overwrites a config.

\begin{lstlisting}[language=ATEShell]
ate setup
sudo ate setup                          # Linux, so the ownership and group work
ate setup --root /tmp/station --no-permissions   # a throwaway Station, for checking
\end{lstlisting}

\section{What it creates}

\texttt{ate setup} makes the data directory with \texttt{reports}, \texttt{logs} and
\texttt{definitions} inside it, and the config directory. On Windows those two are the same
directory; on Linux they are \texttt{/var/lib/ate} and \texttt{/etc/ate}. It then writes a default
\texttt{config.json} --- only if there is no config file there, never over one --- and runs the
same database migration \texttt{ate db migrate} runs (chapter~\ref{ch:database}).

The default config it writes has every key filled in, a fresh \texttt{station.uuid}, the computer's
hostname as \texttt{station.name}, and the paths already written out as real paths rather than as
\texttt{<data>} placeholders, so the Station works wherever setup put it. \texttt{definitions} is
created because the default config names it as a Test Definition directory, and a Test Definition
directory that does not exist is a warning on the Run screen at every scan; setup leaves the
Station quiet until a Media Release is dropped in.

What it prints is the list of directories it created, whether it wrote the config or kept the one
that was there, and the database it migrated:

\begin{lstlisting}[language=ATEShell]
database to migrate: /tmp/station/data/ate.db (database.path in /tmp/station/config/config.json)
created or upgraded database at /tmp/station/data/ate.db (schema version 5)
ate setup:
  created: /tmp/station/data
  created: /tmp/station/data/reports
  created: /tmp/station/data/logs
  created: /tmp/station/data/definitions
  created: /tmp/station/config
  wrote default config: /tmp/station/config/config.json
  database: /tmp/station/data/ate.db
\end{lstlisting}

The database setup migrates is the one the config names, not \texttt{<data>/ate.db} by assumption: it
reads \texttt{config.json} and takes \texttt{database.path} from it, and prints both the path and the
reason it chose it. A Station whose config points the database somewhere else is migrated where the
config says, which is what an installer needs to do (chapter~\ref{ch:upgrades}). If no config can be
read --- it does not exist yet, or it is admin-only and setup is not elevated --- setup falls back to
\texttt{<data>/ate.db} and says so in the same line. Setup never migrates a database it did not print.

\texttt{--root DIR} puts the Station under \texttt{DIR} (\texttt{DIR/data} and
\texttt{DIR/config}); \texttt{--data-dir} and \texttt{--config-dir} override one of them. Those
three exist for packagers and for tests: an ordinary Station uses the platform defaults.

\section{Moving or cloning a Station}

A Station is nothing but its config file and its data directory, so moving one is copying both.
Setup writes every path in the config out in full, not as \texttt{<data>}, so a config copied to a
new machine or a new directory still names the old one: edit \texttt{database.path},
\texttt{reports.directory}, \texttt{log.directory} and \texttt{testDefinitions.directories} to the
new locations, or run \texttt{ate setup --root DIR} there and copy the database into what it made.
Setup never overwrites a config that exists, so it will not rewrite paths for you.

Cloning needs one extra edit. \texttt{station.uuid} is the Station's identity and it is recorded on
every Test Run, and setup writes a config with a fresh one only when the file does not exist. A
copied config carries the original uuid with it, and ATE never notices two Stations claiming the
same one. Blank the value --- \texttt{"uuid": ""} --- on the copy, and the next launch generates and
writes back a new uuid; deleting the key does the same. Two Stations on one database file is a
separate mistake: nothing in the code prevents it, and the runs from both are indistinguishable
except by that uuid.

\section{Permissions, and who can run what}

The privileged part of setup is the ownership and access rights, and it only happens when the
process is elevated --- \texttt{ate setup} run as root on Linux, or as an administrator on
Windows. An unprivileged run prints

\begin{lstlisting}[language=ATEShell]
ate setup: running unprivileged; skipping group/ownership/ACL setup (packagers run this elevated)
\end{lstlisting}

and carries on. It does not fail: the directories and the database are still created, and
\texttt{sudo ate setup} finishes the job later. \texttt{--no-permissions} skips that treatment
without the warning being about privilege.

On Linux, setup creates a system group called \texttt{ate} if it is missing, makes everything under
the data directory group-owned and group-writable, recursively, and sets the setgid bit on the
directories so new files --- \texttt{ate.db-wal}, \texttt{ate.db-shm}, a fresh Test Report ---
inherit the group. \texttt{config.json} and the config directory itself are made root-owned:
the file 0644, the directory 0755. Adding an Operator's account to the \texttt{ate} group is a
step you do yourself; setup creates the group but never adds anyone to it.

On Windows, \texttt{icacls} gives the Users group Modify on the data directory with inherit
flags, so the database, its WAL files and every report and log under it are writable by ordinary
users, and gives \texttt{config.json} an explicit access list: Administrators and SYSTEM full
control, Users read only.

The suite proves this with a recording double that captures every call setup makes, rather than
running \texttt{groupadd}, \texttt{chown} or \texttt{icacls} on the machine the tests run on; the
real Linux commands are exercised in a \texttt{debian:stable} container, gated behind
\texttt{ATE\_DOCKER\_TESTS=1}.

\section{Checking it worked}

\begin{lstlisting}[language=ATEShell]
ate --version
ate db status
\end{lstlisting}

\texttt{ate --version} prints the version the build stamped in (a git tag, or \texttt{0.0.0-dev}
from a checkout). \texttt{ate db status} prints the database's schema version and exits 0; a
missing or outdated database makes it print the refusal and exit 1. If the Station's config is not
the platform default, name it with \texttt{--config}, either before the subcommand
(\texttt{ate --config /tmp/station/config/config.json db status}) or after it
(\texttt{ate db status --config /tmp/station/config/config.json}).

\cando{prepare a Station from nothing, read what setup said, and tell which rights an Operator
needs in order to run a Test Run.}
